\documentclass[10pt,landscape]{article}
\usepackage[landscape,margin=0.3in]{geometry}
\usepackage{multicol}
\usepackage{amsmath}
\usepackage{amssymb}
\usepackage{enumitem}
\usepackage{xcolor}
\usepackage[compact]{titlesec}

% Reduce spacing
\setlength{\parindent}{0pt}
\setlength{\parskip}{1pt}
\setlength{\columnsep}{10pt}

% Section formatting
\titleformat{\section}{\normalsize\bfseries\color{blue!70!black}}{\thesection}{0.5em}{}
\titlespacing*{\section}{0pt}{3pt}{1pt}
\titleformat{\subsection}{\small\bfseries\color{blue!50!black}}{\thesubsection}{0.5em}{}
\titlespacing*{\subsection}{0pt}{2pt}{1pt}

\begin{document}
\begin{multicols}{3}
\small

\begin{center}
\textbf{\large ORF309 Midterm 1 - Clear Cheat Sheet}
\end{center}

\section{Conditional Probability (Key for Every Exam!)}

\subsection{The Basic Formula}
\textbf{Definition:} Probability of $A$ happening, given we know $B$ happened:
$$P(A|B) = \frac{P(A \cap B)}{P(B)}$$

\textbf{What this means:} We "zoom in" to the world where $B$ happened, then ask what fraction of that world contains $A$.

\textbf{Rearranged (multiplication rule):}
$$P(A \cap B) = P(A|B) \cdot P(B)$$

\subsection{Chain Rule (Multiple Events)}
For several events happening together:
$$P(A_1 \cap A_2 \cap A_3) = P(A_1) \cdot P(A_2|A_1) \cdot P(A_3|A_1 \cap A_2)$$

\textbf{Process:} Start with first event, then condition each next event on all previous ones.

\section{Law of Total Probability (LOTP)}

\subsection{When to Use It}
Use LOTP when you want to find $P(A)$ but the path to $A$ goes through different "cases" or "stages".

\subsection{The Formula}
If $B_1, B_2, \ldots, B_n$ partition the sample space (disjoint and cover everything):
$$P(A) = P(A|B_1)P(B_1) + P(A|B_2)P(B_2) + \cdots + P(A|B_n)P(B_n)$$

\textbf{Two-case version (most common):}
$$P(A) = P(A|B)P(B) + P(A|B^c)P(B^c)$$

\subsection{Example: Defective Widgets}
Factory has Line A (produces 60\% of widgets, 2\% defective) and Line B (produces 40\%, 5\% defective). What's probability a random widget is defective?

\textbf{Solution:} Let $D$ = defective, $A$ = from Line A, $B$ = from Line B

$$P(D) = P(D|A)P(A) + P(D|B)P(B)$$
$$= 0.02 \times 0.6 + 0.05 \times 0.4 = 0.012 + 0.020 = 0.032$$

\section{Bayes' Theorem (Updating Beliefs)}

\subsection{The Formula}
$$P(B_j|A) = \frac{P(A|B_j)P(B_j)}{P(A)} = \frac{P(A|B_j)P(B_j)}{\sum_i P(A|B_i)P(B_i)}$$

\textbf{In words:}
\begin{itemize}[leftmargin=*,noitemsep,topsep=0pt]
\item $P(B_j)$ = prior (what we thought before seeing data $A$)
\item $P(A|B_j)$ = likelihood (how likely is data $A$ if $B_j$ is true)
\item $P(B_j|A)$ = posterior (updated belief after seeing data $A$)
\end{itemize}

\subsection{Example: Defective Widget Origin}
Same factory. Given a widget is defective, what's probability it came from Line A?

$$P(A|D) = \frac{P(D|A)P(A)}{P(D)} = \frac{0.02 \times 0.6}{0.032} = \frac{0.012}{0.032} = 0.375$$

\textbf{Interpretation:} Even though Line A produces 60\% of widgets, only 37.5\% of defective widgets come from it (because Line B has higher defect rate).

\section{Independence}

\subsection{What Independence Means}
$A$ and $B$ are independent if knowing one tells you nothing about the other:
$$P(A \cap B) = P(A) \cdot P(B)$$

\textbf{Equivalently:} $P(A|B) = P(A)$ (when $P(B) > 0$)

\subsection{Critical Warning!}
\textbf{Mutually Exclusive $\neq$ Independent!}

If $A \cap B = \emptyset$ (disjoint) and $P(A), P(B) > 0$:
\begin{itemize}[leftmargin=*,noitemsep,topsep=0pt]
\item Then $A$ and $B$ are \textbf{DEPENDENT}
\item Why? If $A$ happens, you know for sure $B$ didn't happen
\end{itemize}

\textbf{Example:} Rolling a die. Let $A$ = "roll 1 or 2", $B$ = "roll 3 or 4". These are disjoint but dependent (knowing one happened tells you the other didn't).

\subsection{Properties of Independence}
If $A \perp B$ (independent), then also:
\begin{itemize}[leftmargin=*,noitemsep,topsep=0pt]
\item $A \perp B^c$ (independent of complement)
\item $A^c \perp B$
\item $A^c \perp B^c$
\end{itemize}

\section{Discrete Distributions (Know These Cold!)}

\subsection{Geometric Distribution: Geom(p)}
\textbf{Models:} Number of trials until first success (when each trial succeeds with probability $p$)

\textbf{PMF:} $P(X = k) = (1-p)^{k-1} \cdot p$ for $k = 1, 2, 3, \ldots$

\textbf{Key formulas:}
\begin{itemize}[leftmargin=*,noitemsep,topsep=0pt]
\item $E[X] = \frac{1}{p}$ (average trials until success)
\item $P(X > k) = (1-p)^k$ (haven't succeeded yet after $k$ tries)
\item $\text{Var}(X) = \frac{1-p}{p^2}$
\end{itemize}

\textbf{MEMORYLESS PROPERTY} (unique to Geometric!):
$$P(X > m + n | X > m) = P(X > n)$$

\textbf{What this means:} If you've already waited $m$ trials without success, the probability of waiting at least $n$ more is the same as if you were starting fresh.

\textbf{Example:} Bus comes with probability 0.1 each minute. You've waited 10 minutes. What's probability you wait at least 20 more?

$$P(X > 30 | X > 10) = P(X > 20) = (0.9)^{20}$$

\subsection{Binomial Distribution: Bin(n, p)}
\textbf{Models:} Number of successes in $n$ independent trials (each succeeds with probability $p$)

\textbf{PMF:} $P(X = k) = \binom{n}{k} p^k (1-p)^{n-k}$ for $k = 0, 1, \ldots, n$

\textbf{Key formulas:}
\begin{itemize}[leftmargin=*,noitemsep,topsep=0pt]
\item $E[X] = np$
\item $\text{Var}(X) = np(1-p)$
\end{itemize}

\textbf{Example:} Flip fair coin 10 times. Number of heads is Bin(10, 0.5):
\begin{itemize}[leftmargin=*,noitemsep,topsep=0pt]
\item $E[\text{heads}] = 10 \times 0.5 = 5$
\item $P(\text{exactly 7 heads}) = \binom{10}{7}(0.5)^7(0.5)^3 = \binom{10}{7}(0.5)^{10}$
\end{itemize}

\subsection{Poisson Distribution: Pois($\lambda$)}
\textbf{Models:} Number of events in fixed time/space when events happen at rate $\lambda$

\textbf{PMF:} $P(N = k) = e^{-\lambda} \frac{\lambda^k}{k!}$ for $k = 0, 1, 2, \ldots$

\textbf{Key formulas:}
\begin{itemize}[leftmargin=*,noitemsep,topsep=0pt]
\item $E[N] = \lambda$
\item $\text{Var}(N) = \lambda$
\end{itemize}

\textbf{When to use:} "Rare events" - when $n$ is large, $p$ is small, and $np = \lambda$ is moderate:
$$\text{Bin}(n, p) \approx \text{Pois}(np)$$

\section{Expectation (The Average)}

\subsection{Basic Definition}
For discrete random variable $X$:
$$E[X] = \sum_{\text{all } x} x \cdot P(X = x)$$

\textbf{What this means:} Weighted average of all possible values, weighted by their probabilities.

\subsection{Linearity of Expectation (THE MOST USEFUL PROPERTY!)}
\textbf{For any random variables $X$ and $Y$ (even dependent!):}
$$E[aX + bY + c] = aE[X] + bE[Y] + c$$

\textbf{Why this is amazing:} No independence needed! Works always!

\textbf{Example:} Roll two dice. Expected sum?

Let $X_1$ = first die, $X_2$ = second die. Even though they're not independent of their sum, we can use linearity:
$$E[X_1 + X_2] = E[X_1] + E[X_2] = 3.5 + 3.5 = 7$$

\subsection{Indicator Trick (Simplify Everything!)}
\textbf{Key idea:} Define indicator random variables for events:
$$1_A = \begin{cases} 1 & \text{if event } A \text{ happens} \\ 0 & \text{otherwise} \end{cases}$$

\textbf{Magic property:} $E[1_A] = P(A)$

\textbf{How to use it:}
\begin{enumerate}[leftmargin=*,noitemsep,topsep=0pt]
\item Write $X$ as sum of indicators: $X = \sum_i 1_{A_i}$
\item Use linearity: $E[X] = \sum_i E[1_{A_i}] = \sum_i P(A_i)$
\end{enumerate}

\textbf{Example:} Bob's exam has 2 questions. He gets question 1 right with prob 0.7, question 2 right with prob 0.6. Each worth 10 points. Expected score?

Let $X$ = score = $10 \cdot 1_{Q1} + 10 \cdot 1_{Q2}$

$$E[X] = 10 \cdot P(Q1) + 10 \cdot P(Q2) = 10(0.7) + 10(0.6) = 13$$

\textbf{Note:} Works even if questions aren't independent!

\subsection{Conditional Expectation}
\textbf{Definition:} Expected value of $X$ given event $B$ happened:
$$E[X | B] = \sum_{\text{all } x} x \cdot P(X = x | B)$$

\textbf{Law of Total Expectation (Tower Property):}
$$E[X] = \sum_i E[X | B_i] P(B_i)$$

where $B_1, \ldots, B_n$ partition the sample space.

\textbf{When to use:} When computing $E[X]$ directly is hard, but $E[X | \text{something}]$ is easier.

\textbf{Example:} Expected widget defects. If from Line A: $E[\text{defects}|A] = 0.02$, if from Line B: $E[\text{defects}|B] = 0.05$. Line A makes 60\% of widgets.

$$E[\text{defects}] = E[\text{defects}|A] \cdot P(A) + E[\text{defects}|B] \cdot P(B)$$
$$= 0.02 \times 0.6 + 0.05 \times 0.4 = 0.032$$

\section{Variance (Measuring Spread)}

\subsection{Definition \& Key Formula}
$$\text{Var}(X) = E[(X - \mu)^2] = E[X^2] - (E[X])^2$$

\textbf{Second formula is usually easier!} Calculate $E[X^2]$ and $E[X]$, then:
$$\text{Var}(X) = E[X^2] - (E[X])^2$$

\subsection{Properties}
\begin{itemize}[leftmargin=*,noitemsep,topsep=0pt]
\item $\text{Var}(aX + b) = a^2 \text{Var}(X)$ (adding constant doesn't change spread)
\item $\text{Var}(X) \geq 0$ always
\item $\text{Var}(X) = 0$ only if $X$ is constant
\end{itemize}

\subsection{Variance of Sum (NEEDS Independence!)}
\textbf{If $X$ and $Y$ are independent:}
$$\text{Var}(X + Y) = \text{Var}(X) + \text{Var}(Y)$$

\textbf{Critical:} Unlike expectation, this \textbf{requires independence}!

\textbf{Example:} Bob's exam (continued). Variance of score?

Assuming questions independent:
$$\text{Var}(X) = \text{Var}(10 \cdot 1_{Q1}) + \text{Var}(10 \cdot 1_{Q2})$$
$$= 100 \cdot \text{Var}(1_{Q1}) + 100 \cdot \text{Var}(1_{Q2})$$

For Bernoulli with $p$: $\text{Var} = p(1-p)$

$$= 100[0.7(0.3)] + 100[0.6(0.4)] = 21 + 24 = 45$$

\section{Problem-Solving Strategies}

\subsection{When to Use Each Technique}

\textbf{Use LOTP when:}
\begin{itemize}[leftmargin=*,noitemsep,topsep=0pt]
\item Problem has multiple stages or cases
\item You know conditional probabilities but not direct probability
\item Keywords: "first..., then...", "depending on..."
\end{itemize}

\textbf{Use Bayes when:}
\begin{itemize}[leftmargin=*,noitemsep,topsep=0pt]
\item You want to "reverse" a conditional probability
\item You know $P(\text{effect}|\text{cause})$ but want $P(\text{cause}|\text{effect})$
\item Keywords: "given that..., what caused it?"
\end{itemize}

\textbf{Use indicator trick when:}
\begin{itemize}[leftmargin=*,noitemsep,topsep=0pt]
\item Computing $E[X]$ where $X$ counts something
\item Can break $X$ into sum of events
\item Especially good when events might be dependent
\end{itemize}

\textbf{Use geometric distribution when:}
\begin{itemize}[leftmargin=*,noitemsep,topsep=0pt]
\item "Waiting time" until first success
\item Repeated independent trials
\item Keywords: "how many tries until...", "first time..."
\end{itemize}

\subsection{Step-by-Step for Exam Problems}

\textbf{For probability problems:}
\begin{enumerate}[leftmargin=*,noitemsep,topsep=0pt]
\item Define events clearly (write out what $A$, $B$, etc. mean)
\item Identify if events are independent (check if one affects the other)
\item Look for partitions (use LOTP)
\item Check if you need to reverse conditional (use Bayes)
\end{enumerate}

\textbf{For expectation problems:}
\begin{enumerate}[leftmargin=*,noitemsep,topsep=0pt]
\item Can you write $X$ as sum of indicators? (Do it!)
\item Is this a named distribution? (Use formula)
\item Can you condition on something simpler? (Use tower property)
\item Use linearity freely (no independence needed)
\end{enumerate}

\textbf{For variance problems:}
\begin{enumerate}[leftmargin=*,noitemsep,topsep=0pt]
\item First find $E[X]$
\item Then find $E[X^2]$ (often via sum of indicators squared)
\item Use $\text{Var}(X) = E[X^2] - (E[X])^2$
\item Check if variables are independent before adding variances!
\end{enumerate}

\section{Common Exam Mistakes to Avoid}

\begin{enumerate}[leftmargin=*,noitemsep,topsep=0pt]
\item \textbf{Thinking disjoint = independent.} NO! Disjoint events are dependent (if one happens, other can't).

\item \textbf{Forgetting linearity always works for E[·].} You can split expectations even when variables are dependent!

\item \textbf{Adding variances without independence.} $\text{Var}(X+Y) = \text{Var}(X) + \text{Var}(Y)$ ONLY when independent.

\item \textbf{Not recognizing geometric distribution.} Keywords: "first success", "how many tries", "until..." $\Rightarrow$ probably Geometric.

\item \textbf{Forgetting memoryless property.} If problem says "already waited $m$ trials" and asks about more waiting $\Rightarrow$ use memoryless: $P(X > m+n | X > m) = P(X > n) = (1-p)^n$.

\item \textbf{Confusing $P(A|B)$ with $P(B|A)$.} These are different! To reverse, use Bayes.

\item \textbf{Not using indicator trick.} If you're counting events, write as sum of indicators: $X = \sum 1_{A_i}$, then $E[X] = \sum P(A_i)$.

\item \textbf{Forgetting $E[X^2] \neq (E[X])^2$.} These are only equal when $\text{Var}(X) = 0$.

\item \textbf{Computing expectations "the hard way".} Before computing $\sum x \cdot P(X=x)$, ask: Is this a named distribution? Can I use indicators? Can I condition on something?
\end{enumerate}

\section{Quick Reference Formulas}

\subsection{Must-Know Distributions}
\begin{itemize}[leftmargin=*,noitemsep,topsep=0pt]
\item Geom$(p)$: $E = 1/p$, $P(X>k) = (1-p)^k$, memoryless
\item Bin$(n,p)$: $E = np$, Var $= np(1-p)$
\item Pois$(\lambda)$: $E = \lambda$, Var $= \lambda$
\item Bernoulli$(p)$: $E = p$, Var $= p(1-p)$
\end{itemize}

\subsection{Key Identities}
\begin{itemize}[leftmargin=*,noitemsep,topsep=0pt]
\item $P(A \cup B) = P(A) + P(B) - P(A \cap B)$
\item $P(A^c) = 1 - P(A)$
\item $P(A \cap B) = P(A|B)P(B)$
\item $E[1_A] = P(A)$ (indicator trick)
\item $\text{Var}(X) = E[X^2] - (E[X])^2$
\item Independence: $X \perp Y \Rightarrow E[XY] = E[X]E[Y]$
\end{itemize}

\end{multicols}
\end{document}
